\documentclass[10pt,a4paper]{amsart}
\usepackage{amsmath,amssymb,amsthm,graphicx}
% Retained minimal preamble for the Stacks Project chapter flat.tex (More on Flatness).
% Verbatim excerpt of the chapter preamble: Xy-pic setup (native preamble lines 14-18) and the native
% theorem declarations with the mathematical macros (native preamble lines 154-217). Copied mechanically;
% nothing added, removed or redefined. The wrapper class and the amsmath/amssymb/amsthm/graphicx base
% packages are supplied by the assembler, so the class/externaldocument part of the native preamble is omitted.
% For commutative diagrams we use Xy-pic
\usepackage[all]{xy}

% We use 2cell for 2-commutative diagrams.
\xyoption{2cell}
% Theorem environments.
%
\theoremstyle{plain}
\newtheorem{theorem}[subsection]{Theorem}
\newtheorem{proposition}[subsection]{Proposition}
\newtheorem{lemma}[subsection]{Lemma}

\theoremstyle{definition}
\newtheorem{definition}[subsection]{Definition}
\newtheorem{example}[subsection]{Example}
\newtheorem{exercise}[subsection]{Exercise}
\newtheorem{situation}[subsection]{Situation}

\theoremstyle{remark}
\newtheorem{remark}[subsection]{Remark}
\newtheorem{remarks}[subsection]{Remarks}

\numberwithin{equation}{subsection}

% Macros
%
\def\lim{\mathop{\mathrm{lim}}\nolimits}
\def\colim{\mathop{\mathrm{colim}}\nolimits}
\def\Spec{\mathop{\mathrm{Spec}}}
\def\Hom{\mathop{\mathrm{Hom}}\nolimits}
\def\Ext{\mathop{\mathrm{Ext}}\nolimits}
\def\SheafHom{\mathop{\mathcal{H}\!\mathit{om}}\nolimits}
\def\SheafExt{\mathop{\mathcal{E}\!\mathit{xt}}\nolimits}
\def\Sch{\mathit{Sch}}
\def\Mor{\mathop{\mathrm{Mor}}\nolimits}
\def\Ob{\mathop{\mathrm{Ob}}\nolimits}
\def\Sh{\mathop{\mathit{Sh}}\nolimits}
\def\NL{\mathop{N\!L}\nolimits}
\def\CH{\mathop{\mathrm{CH}}\nolimits}
\def\proetale{{pro\text{-}\acute{e}tale}}
\def\etale{{\acute{e}tale}}
\def\QCoh{\mathit{QCoh}}
\def\Ker{\mathop{\mathrm{Ker}}}
\def\Im{\mathop{\mathrm{Im}}}
\def\Coker{\mathop{\mathrm{Coker}}}
\def\Coim{\mathop{\mathrm{Coim}}}

% Boxtimes
%
\DeclareMathSymbol{\boxtimes}{\mathbin}{AMSa}{"02}

%
% Macros for moduli stacks/spaces
%
\def\QCohstack{\mathcal{QC}\!\mathit{oh}}
\def\Cohstack{\mathcal{C}\!\mathit{oh}}
\def\Spacesstack{\mathcal{S}\!\mathit{paces}}
\def\Quotfunctor{\mathrm{Quot}}
\def\Hilbfunctor{\mathrm{Hilb}}
\def\Curvesstack{\mathcal{C}\!\mathit{urves}}
\def\Polarizedstack{\mathcal{P}\!\mathit{olarized}}
\def\Complexesstack{\mathcal{C}\!\mathit{omplexes}}
% \Pic is the operator that assigns to X its picard group, usage \Pic(X)
% \Picardstack_{X/B} denotes the Picard stack of X over B
% \Picardfunctor_{X/B} denotes the Picard functor of X over B
\def\Pic{\mathop{\mathrm{Pic}}\nolimits}
\def\Picardstack{\mathcal{P}\!\mathit{ic}}
\def\Picardfunctor{\mathrm{Pic}}
\def\Deformationcategory{\mathcal{D}\!\mathit{ef}}

% --- appended 2026-10-12 (Chief): the project own definitions of the reference/slogan/history
% display environments, copied mechanically from the shared Stacks preamble lines 8-12
% (\usepackage{verbatim} plus the three \newenvironment lines, which the chapter preamble keeps
% verbatim from the project preamble). Nothing added, removed or redefined: these environments
% carry no mathematical content and the project itself typesets them as comments. Without these
% three lines a counted statement whose own text uses them fails with Environment slogan undefined.
\usepackage{verbatim}
\newenvironment{reference}{\comment}{\endcomment}
\newenvironment{slogan}{\comment}{\endcomment}
\newenvironment{history}{\comment}{\endcomment}
\usepackage{cleveref}
\crefname{originalreference}{original reference}{original references}
\Crefname{originalreference}{Original reference}{Original references}
\makeatletter
\newlabel{situation-complex-and-divisor}{{0F8S}{source}{}{sourcepointer-1}{}}
\newlabel{situation-complex-and-divisor@cref}{{[originalreference][0][]0F8S}{[1][0][]source}{}{}{}}
\makeatother
\crefname{originallemma}{original lemma}{original lemmas}
\Crefname{originallemma}{Original lemma}{Original lemmas}
\makeatletter
\newlabel{lemma-complex-and-divisor-blowup-pre}{{0GTY}{source}{}{sourcepointer-2}{}}
\newlabel{lemma-complex-and-divisor-blowup-pre@cref}{{[originallemma][0][]0GTY}{[1][0][]source}{}{}{}}
\makeatother
\crefname{originallemma}{original lemma}{original lemmas}
\Crefname{originallemma}{Original lemma}{Original lemmas}
\makeatletter
\newlabel{lemma-complex-and-divisor-blowup}{{0F8U}{source}{}{sourcepointer-3}{}}
\newlabel{lemma-complex-and-divisor-blowup@cref}{{[originallemma][0][]0F8U}{[1][0][]source}{}{}{}}
\makeatother
\crefname{originallemma}{original lemma}{original lemmas}
\Crefname{originallemma}{Original lemma}{Original lemmas}
\makeatletter
\newlabel{lemma-pullback-triple-ideals-good}{{0F8T}{source}{}{sourcepointer-4}{}}
\newlabel{lemma-pullback-triple-ideals-good@cref}{{[originallemma][0][]0F8T}{[1][0][]source}{}{}{}}
\makeatother
\crefname{originallemma}{original lemma}{original lemmas}
\Crefname{originallemma}{Original lemma}{Original lemmas}
\makeatletter
\newlabel{lemma-complex-and-divisor-eta-pull}{{0F9X}{source}{}{sourcepointer-5}{}}
\newlabel{lemma-complex-and-divisor-eta-pull@cref}{{[originallemma][0][]0F9X}{[1][0][]source}{}{}{}}
\makeatother
\title{The blowup making a complex and divisor into a good triple is compatible with base change, where the new scheme is the strict transform}
\author{The Stacks Project Authors}
\date{stacks-project:flat.tex:local-snapshot-20261011; GNU FDL 1.2 or later}
\begin{document}
\maketitle
\noindent\textit{Editorial note.} All delivered mathematical text is the author's own, copied from the pinned Stacks Project chapter file flat.tex; no statement or proof step was rewritten, reordered or removed.\par\medskip
\noindent\textit{Editorial note.} The preamble is the chapter's own retained minimal preamble (sources/stacks-flat/retained-preamble.tex), a verbatim excerpt of the native Xy-pic setup and of the native theorem declarations and mathematical macros, followed by the project own definitions of the reference, slogan and history display environments (verbatim package plus three newenvironment lines) that the retained file carries at its end and that the project itself typesets as comments. No macro was invented or redefined, and the wrapper class and the amsmath/amssymb/amsthm/graphicx base packages are supplied by the assembler.\par\medskip
\noindent\textit{Editorial note.} The counted statement and proof contain no cross reference, no citation and no included file, so the delivered body is self contained.\par\medskip
\noindent\textit{Editorial note.} Canonical location of the counted proposition: Stacks tag 0F8V (sources/stacks-flat/pinned-original-source.tex line 13731).\par\medskip
\noindent\textit{Allowed standard prerequisites (editorial).}\begin{itemize}
\item The blowup of a complex and divisor data set and its universal property. Same chapter, Lemma labelled lemma-complex-and-divisor-blowup-pre, cited at proof lines 13749-13756 and in the statement. Produces the unique morphism h' : Y' -> X' compatible with b and c.
\item Affine charts and the pullback of the ideals I\_i(-,-). Same chapter, Lemma labelled lemma-pullback-triple-ideals-good, cited at proof lines 13770-13775, together with the chart description from the proof of Lemma labelled lemma-complex-and-divisor-blowup-pre cited at proof lines 13764-13768. Identifies the restriction of c with the blowup in the pullback ideal, which is the strict transform.
\item Pullback of the derived twist eta along the base change. Same chapter, Lemma labelled lemma-complex-and-divisor-eta-pull, cited at proof lines 13786-13788. Gives L eta\_\{J'\} L(h o c)\textasciicircum{}*M = L(Y' -> X')\textasciicircum{}*Q.
\item The complex Q attached to the blowup. Same chapter, Lemma labelled lemma-complex-and-divisor-blowup, invoked in the statement. Fixes Q = L eta\_\{I'\} Lb\textasciicircum{}*M on X'.
\end{itemize}\medskip
\ifdefined\hypertarget\hypertarget{sourcepointer-1}{}\fi
\noindent\textit{Source reference.} Original reference 0F8S: sources/stacks-flat/pinned-original-source.tex, source lines 13460--13460. Same-chapter pointer into the author environment situation at source line 13460; invoked at proof line None of the counted proof. It is retained only as a navigation pointer; every load-bearing step is inside the delivered span.\par\medskip
\ifdefined\hypertarget\hypertarget{sourcepointer-2}{}\fi
\noindent\textit{Source reference.} Original lemma 0GTY: sources/stacks-flat/pinned-original-source.tex, source lines 13588--13588. Same-chapter pointer into the author environment lemma at source line 13588; invoked at proof line 13752 of the counted proof. It is retained only as a navigation pointer; every load-bearing step is inside the delivered span.\par\medskip
\ifdefined\hypertarget\hypertarget{sourcepointer-3}{}\fi
\noindent\textit{Source reference.} Original lemma 0F8U: sources/stacks-flat/pinned-original-source.tex, source lines 13717--13717. Same-chapter pointer into the author environment lemma at source line 13717; invoked at proof line 13767 of the counted proof. It is retained only as a navigation pointer; every load-bearing step is inside the delivered span.\par\medskip
\ifdefined\hypertarget\hypertarget{sourcepointer-4}{}\fi
\noindent\textit{Source reference.} Original lemma 0F8T: sources/stacks-flat/pinned-original-source.tex, source lines 13488--13488. Same-chapter pointer into the author environment lemma at source line 13488; invoked at proof line 13773 of the counted proof. It is retained only as a navigation pointer; every load-bearing step is inside the delivered span.\par\medskip
\ifdefined\hypertarget\hypertarget{sourcepointer-5}{}\fi
\noindent\textit{Source reference.} Original lemma 0F9X: sources/stacks-flat/pinned-original-source.tex, source lines 13570--13570. Same-chapter pointer into the author environment lemma at source line 13570; invoked at proof line 13786 of the counted proof. It is retained only as a navigation pointer; every load-bearing step is inside the delivered span.\par\medskip
\begin{lemma}
\label{lemma-complex-and-divisor-blowup-base-change}
In Situation \ref{situation-complex-and-divisor} let $h : Y \to X$
be a morphism of schemes such that the pullback $E = h^{-1}D$
is defined. Let $b : X' \to X$, resp.\ $c : Y' \to Y$ be as constructed in
Lemma \ref{lemma-complex-and-divisor-blowup-pre} for
$D \subset X$ and $M$, resp.\ $E \subset Y$ and $Lh^*M$.
Then $Y'$ is the strict transform of $Y$ with respect to $b : X' \to X$
(see proof for a precise formulation of this) and
$$
L\eta_{\mathcal{J}'}L(h \circ c)^*M = L(Y' \to X')^*Q
$$
where $Q = L\eta_{\mathcal{I}'}Lb^*M$ as in
Lemma \ref{lemma-complex-and-divisor-blowup}.
In particular, if $(Y, E, Lh^*M)$ is a good triple and
$k : Y \to X'$ is the unique morphism such that
$h = b \circ k$, then $L\eta_\mathcal{J}Lh^*M = Lk^*Q$.
\end{lemma}

\begin{proof}
Denote $E' = c^{-1}E$. Then $(Y', E', L(h \circ c)^*M)$ is a good
triple. Hence by the universal property of
Lemma \ref{lemma-complex-and-divisor-blowup-pre}
there is a unique morphism
$$
h' : Y' \longrightarrow X'
$$
such that $b \circ h' = h \circ c$. In particular, there is a morphism
$(h', c) : Y' \to X' \times_X Y$. We claim that given $W \subset X$
quasi-compact open, such that $b^{-1}(W) \to W$ is a blowing up,
this morphism identifies $Y'|_W$ with the strict transform of $Y_W$
with respect to $b^{-1}(W) \to W$. In turn, to see this is true
is a local question on $W$, and we may therefore prove the
statement over an affine chart. We do this in the next paragraph.

\medskip\noindent
Let $(U, A, f, M^\bullet)$ be an affine chart for $(X, D, M)$.
Recall from the proof of Lemma \ref{lemma-complex-and-divisor-blowup}
that the restriction of $b : X' \to X$ to $U$ is the blowing
up of $U = \Spec(A)$ in the product of the ideals $I_i(M^\bullet, f)$.
Now if $V = \Spec(B) \subset Y$ is any affine open with $h(V) \subset U$,
then $(V, B, g, M^\bullet \otimes_A B)$ is an affine chart for
$(Y, E, Lh^*M)$ where $g \in B$ is the image of $f$, see
Lemma \ref{lemma-pullback-triple-ideals-good}.
Hence the restriction of $c : Y' \to Y$ to $V$
is the blowing up in the product of the ideals
$I_i(M^\bullet, f)B$, i.e., the morphism $c : Y' \to Y$ over
$h^{-1}(U)$ is the blowing up of $h^{-1}(U)$ in the
ideal $\prod I_i(M^\bullet, f) \mathcal{O}_{h^{-1}(U)}$.
Since this is also true for the strict transform, we see that
our claim on strict transforms is true.

\medskip\noindent
Having said this the equality
$L\eta_{\mathcal{J}'}L(h \circ c)^*M = L(Y' \to X')^*Q$
follows from
Lemma \ref{lemma-complex-and-divisor-eta-pull}.
The final statement is a special case of this
(namely, the case where $c = \text{id}_Y$ and $k = h'$).
\end{proof}

\end{document}
